\input{sol_pre.tex}
\input{q_2025-05-25_1.tex}
\begin{document}
\input{sol_head.tex}

\Savol{1}
$49^5\cdot12^3=7^{10}\cdot2^6\cdot3^3$, \ $42^n=2^n\cdot3^n\cdot7^n$ $\Rightarrow$ $n\ge10$.
\Javob{D}

\Savol{2}
$\dfrac{2026!}{2025!+2024!}=\dfrac{2026\cdot2025\cdot2024!}{2024!\,(2025+1)}=2025$, \quad $\dfrac{10!}{9!+8!}=\dfrac{10\cdot9\cdot8!}{8!\,(9+1)}=9$.
\[\sqrt{2025}-\sqrt9=45-3=42.\]
\Javob{B}

\Savol{3}
Har biri $n$ ta detal: \ $\dfrac n{20}-\dfrac n{21}=\dfrac16$ (soat) $\Rightarrow$ $\dfrac n{420}=\dfrac16$ $\Rightarrow$ $n=70$.
\Javob{B}

\Savol{4}
Reja: $4020:1{,}2=3350$. \ Rejadan taşqari: $4020-3350=670$.
\Javob{C}

\Savol{5}
$a^4+4b^4=\left(a^2+2b^2+2ab\right)\left(a^2+2b^2-2ab\right)$; \ $a=1$, $b=2^{24}$: \ $2^{98}+1=1+4\cdot2^{96}$:
\[2^{98}+1=\left(2^{49}+2^{25}+1\right)\left(2^{49}-2^{25}+1\right)\ \Rightarrow\ \text{javob }2^{49}-2^{25}+1.\]
\Javob{B}

\Savol{6}
$a^6=26$, \ $b^6=25$, \ $c^6=27$ $\Rightarrow$ $c>a>b$.
\Javob{C}

\Savol{7}
$x=2024$: \ $(x+1)^3-(x-1)^3=6x^2+2$ $\Rightarrow$ $a-2=6\cdot2024^2$, \ $\sqrt{\dfrac{a-2}6}=2024$.
\Javob{D}

\Savol{8}
$S_m=\dfrac{a_1+a_m}2\cdot m=45m=270$ $\Rightarrow$ $m=6$.
\Javob{B}

\Savol{9}
$S_9=9a_5=369$ $\Rightarrow$ $a_5=41$, $d=10$, $a_9=81$. \ $b_9=q^8=81$ $\Rightarrow$ $q^2=3$, \ $b_7=q^6=27$.
\Javob{A}

\Savol{10}
$625^{-(a^2+4ab)}=625^{\frac{3a^2-10ab}3}$ $\Rightarrow$ $-3a^2-12ab=3a^2-10ab$ $\Rightarrow$ $2a(3a+b)=0$ $\Rightarrow$ $\dfrac ab=-\dfrac13$.
\Javob{B}

\Savol{11}
$(a+b+c)^2=0$ $\Rightarrow$ $ab+bc+ca=-\frac12$; \ $(ab+bc+ca)^2=a^2b^2+b^2c^2+c^2a^2+2abc(a+b+c)$ $\Rightarrow$ $\sum a^2b^2=\frac14$.
\[a^4+b^4+c^4=\left(a^2+b^2+c^2\right)^2-2\sum a^2b^2=1-\frac12=\frac12.\]
\Javob{A}

\Savol{12}
$-\dfrac\pi3+\dfrac\pi6=-\dfrac\pi6$.
\Javob{A}

\Savol{13}
$(\sin\alpha-\cos\alpha)^2=\frac15$ $\Rightarrow$ $\sin2\alpha=\frac45$; \ $\sin\alpha+\cos\alpha=\sqrt{1+\frac45}=\frac3{\sqrt5}$.\\
$\cos2\alpha=(\cos\alpha-\sin\alpha)(\cos\alpha+\sin\alpha)=-\frac1{\sqrt5}\cdot\frac3{\sqrt5}=-\frac35$: \ $\frac45+\frac35=\frac75$.
\Javob{D}

\Savol{14}
$t=222^x>0$: \ $t^3+t=222$ — çap tomon ösuvçi, yagona ildiz ($t=6$ dan sal kiçik). $1$ ta yeçim.
\Javob{B}

\Savol{15}
$\dfrac{\log_{12}\log_8x}{\log_{12}0{,}25}=\log_{0{,}25}\left(\log_8x\right)$ $\Rightarrow$ çap tomon $\log_8x>0$.\\
Öng tomon $\log_{\frac18}3<0$ — yeçim yöq.
\Javob{B}

\Savol{16}
$|x-4|+3=\sqrt{4x+17}$.\\
$x\ge4$: $(x-1)^2=4x+17$ $\Rightarrow$ $x^2-6x-16=0$ $\Rightarrow$ $x=8$.\\
$x<4$: $(7-x)^2=4x+17$ $\Rightarrow$ $x^2-18x+32=0$ $\Rightarrow$ $x=2$ ($16$ — mos emas). \ $8+2=10$.
\Javob{A}

\Savol{17}
$t=\sqrt x\ge0$: \ $t^3-8t^2+19t-12=0$ $\Rightarrow$ $(t-1)(t-3)(t-4)=0$ $\Rightarrow$ $x=1,\,9,\,16$. \ Yiğindi $26$.
\Javob{B}

\Savol{18}
$2x-4>0$ va $(x-2)(x-8)\ge0$ $\Rightarrow$ $x\ge8$. Kvadratga oşiramiz: $x^2-10x+16<4x^2-16x+16$ $\Rightarrow$ $3x(x-2)>0$ — bajariladi.\\
$x\in[8;20]$: $13$ ta butun son.
\Javob{D}

\Savol{19}
$x^2+9<25$ $\Rightarrow$ $x^2<16$: natural $1,2,3$ — $3$ ta.
\Javob{A}

\Savol{20}
$6f(x)-5=x^2$ $\Rightarrow$ $f(2)=\dfrac{4+5}6=\dfrac32$.
\Javob{B}

\Savol{21}
$t=\tg\frac x2$: \ $1+\sin x+\cos x=\dfrac{2+2t}{1+t^2}$, \ $dx=\dfrac{2\,dt}{1+t^2}$; \ $x=\frac\pi2\to t=1$, \ $x=\frac{2\pi}3\to t=\sqrt3$:
\[\int_1^{\sqrt3}\frac{dt}{1+t}=\ln\frac{1+\sqrt3}2.\]
\Javob{A}

\Savol{22}
$g(f)=1+\dfrac1f=\dfrac{x-1}{x-2}=1+\dfrac1{x-2}$ $\Rightarrow$ $f(x)=x-2$, \ $f(2023)=2021$.
\Javob{B}

\Savol{23}
$c=\cos x$: \ $f=\dfrac{(c+3)^2}{1-c^2}$, \ $\dfrac{d f}{dc}=\dfrac{2(c+3)(1+3c)}{\left(1-c^2\right)^2}=0$ $\Rightarrow$ $c=-\dfrac13$ (oraliqda bor, masalan $x\in\left(\frac\pi2;\pi\right)$).
\[f_{\min}=\frac{\left(\frac83\right)^2}{\frac89}=8.\]
\Javob{C}

\Savol{24}
$AD=2r$ — diametr, $AB=r$ ($BC$ aylanaga urinadi): \ $\dfrac{\pi r^2}{2r\cdot r}=\dfrac\pi2$.
\Javob{A}

\Savol{25}
$O$ — bissektrisalar kesişişi: $\angle AOB=90^\circ+\dfrac{\angle C}2=150^\circ$ $\Rightarrow$ $\angle C=120^\circ$.\\
Kosinuslar teoremasi, $AC=b$: \ $169=64+b^2+8b$ $\Rightarrow$ $b^2+8b-105=0$ $\Rightarrow$ $b=7$.
\[S=\frac12\cdot8\cdot7\cdot\sin120^\circ=14\sqrt3.\]
\Javob{C}

\Savol{26}
$F$ — $A$ bissektrisasida: $F$ dan $AD$ va $AB$ gaça masofalar teng; $B$ bissektrisasida: $AB$ va $BC$ gaça masofalar teng.\\
Demak $F$ dan $AD$ va $BC$ gaça masofalar teng, ya'ni $F$ — $AD$ va $BC$ orasidagi örta çiziqda: $h_F=\dfrac h2$.
\[S_{AFD}=\frac12\cdot AD\cdot\frac h2=\frac14S_{ABCD}=12.\]
\Javob{D}

\Savol{27}
Muntazam uçburçak içidagi nuqtadan tomonlargaça masofalar yiğindisi balandlikka teng: $h=\dfrac{\sqrt3}2\cdot4\sqrt3=6$.
\[2d+1=6\ \Rightarrow\ d=2{,}5.\]
\Javob{B}

\Savol{28}
Tomon $s$. Har bir yoy $45^\circ$. \ $\angle FBC$ — $3$ ta yoyga tiraladi: $67{,}5^\circ$; \ $\angle BCM$ — $2$ ta yoyga: $45^\circ$ $\Rightarrow$ $\angle BLC=67{,}5^\circ$, ya'ni $LC=BC=s$.
\[S_1=\frac12\cdot s\cdot s\cdot\sin45^\circ=\frac{\sqrt2}4s^2.\]
$MC=s\left(1+\sqrt2\right)$ $\Rightarrow$ $ML=s\sqrt2$. $BF$ — markazdan ötadi, sakkizburçakni teng ikkiga böladi: yarmi $\left(1+\sqrt2\right)s^2$.\\
$ABLM$ — trapetsiya ($AB\parallel ML$), balandligi $\frac s{\sqrt2}$: \ $S_{ABLM}=\dfrac{s+s\sqrt2}2\cdot\dfrac s{\sqrt2}=\dfrac{2+\sqrt2}4s^2$.
\[S_2=\left(1+\sqrt2\right)s^2-\frac{2+\sqrt2}4s^2=\frac{2+3\sqrt2}4s^2,\qquad \frac{S_2}{S_1}=\frac{2+3\sqrt2}{\sqrt2}=3+\sqrt2.\]
\Javob{C}

\Savol{29}
$C$ va $A$ orasidagi balandlik $s\sqrt3=6$ $\Rightarrow$ $s=2\sqrt3$. $B(0;3)$, $BE=2s=4\sqrt3$ $\Rightarrow$ $E\left(4\sqrt3;\,3\right)$.\\
Köpaytma: $12\sqrt3$.
\Javob{D}

\Savol{30}
Katta burçak $21$ tomon qarşisida. $p=24$, $S=\sqrt{24\cdot14\cdot7\cdot3}=84$, balandlik $h=\dfrac{2\cdot84}{21}=8$.\\
Uç perpendikulyar teoremasi: \ $\sqrt{15^2+8^2}=17$.
\Javob{A}

\Savol{31}
$(A\cup B)\cap C=\{g,h,i,j,k,l\}$ — $6$ ta element: \ $2^6=64$.
\Javob{B}

\Savol{32}
Har xonaga $10$ ta raqam: $10^6$.
\Javob{D}

\newpage
\Savol{33}
Oltiburçak tomoni $2$: qisqa diagonal $2\sqrt3$, katta diagonal $4$, qarama-qarşi tomonlar orasidagi masofa $2\sqrt3$.

\textbf{33.} $AC\parallel F_1D_1$, $AC=F_1D_1=2\sqrt3$; $\angle ACD=90^\circ$ va $AC\perp CC_1$ $\Rightarrow$ $AC\perp CD_1$. Töğri törtburçak: $CD_1=2\sqrt2$:
\[S=2\sqrt3\cdot2\sqrt2=4\sqrt6\ \text{— \textbf{B}}.\]

\textbf{34.} $BE\parallel C_1D_1$: trapetsiya, asoslari $4$ va $2$. $BE$ va $CD$ orasidagi masofa $\sqrt3$, balandlik böyiça $2$ $\Rightarrow$ trapetsiya balandligi $\sqrt{3+4}=\sqrt7$:
\[S=\frac{4+2}2\cdot\sqrt7=3\sqrt7\ \text{— \textbf{E}}.\]

\textbf{35.} $AB\parallel E_1D_1$. Kesim tekisligi $AB$ dan $E_1D_1$ gaça ötadi, örta balandlikda $CF$ katta diagonal sathidan ötadi — kesim oltiburçak, uning asosga proyeksiyasi butun oltiburçak: $S_0=6\cdot\frac{\sqrt3}4\cdot4=6\sqrt3$.\\
Tekisliklar orasidagi burçak: $\tg\varphi=\dfrac2{2\sqrt3}$ $\Rightarrow$ $\varphi=30^\circ$:
\[S=\frac{S_0}{\cos30^\circ}=\frac{6\sqrt3}{\sqrt3/2}=12\ \text{— \textbf{A}}.\]
\Javob{33 — B, \ 34 — E, \ 35 — A}

\newpage
\Savol{36}
$y=a-|x|$ — uçi $(0;a)$ dagi «teskari burçak», sistema $x\to-x$ da özgarmaydi. Har bir tarmoq ($x\ge0$ da $y=a-x$) aylanani kesadi, agar $\frac{|a|}{\sqrt2}\le6$.

$-6<a<6$: uç $(0;a)$ aylana içida — har tarmoq bitta nuqtada: $2$ ta yeçim.\\
$a=6$: uç aylanada; $x^2+(6-x)^2=36$ $\Rightarrow$ $x=0$ yoki $6$: \ $(0;6)$, $(\pm6;0)$ — $3$ ta yeçim.\\
$a=-6$: faqat $(0;-6)$ — $1$ ta. \ $6<a<6\sqrt2$: har tarmoq ikki nuqtada — $4$ ta. \ $a=6\sqrt2$: uriniş — $2$ ta.

a) $a=6$.

b) $6<a<6\sqrt2\approx8{,}49$: butunlari $7$, $8$; \ $7+8=15$.
\Javob{a) 6; \quad b) 15}

\newpage
\Savol{37}
$\cos2x-1=-2\sin^2x$, $s=\sin x$: \ maxraj $-2s^2-s\ne0$ $\Rightarrow$ $s\ne0$, $s\ne-\frac12$.
\[\frac{0{,}5-1+2s^2}{-2s^2-s}=-2s^2+s\ \Rightarrow\ 2s^2-\frac12=4s^4-s^2\ \Rightarrow\ 8s^4-6s^2+1=0\ \Rightarrow\ s^2=\frac12\ \text{yoki}\ \frac14.\]
Yeçimlar: $\sin x=\pm\frac{\sqrt2}2$ $\Rightarrow$ $x=\frac\pi4+\frac{\pi k}2$; \ $\sin x=\frac12$ $\Rightarrow$ $x=\frac\pi6+2\pi k$, $\frac{5\pi}6+2\pi k$ ($\sin x=-\frac12$ çetlatiladi).

a) Eng kiçik musbat $\frac\pi6$, eng katta manfiy $-\frac\pi4$: \ $\frac\pi6-\frac\pi4=-\frac\pi{12}$.

b) $[-\pi;\pi]$: $\pm\frac\pi4$, $\pm\frac{3\pi}4$, $\frac\pi6$, $\frac{5\pi}6$ — $6$ ta.
\Javob{a) $-\frac\pi{12}$; \quad b) 6}

\newpage
\Savol{38}
Uçi $x=2$ da, ildizlar simmetrik: $m=4-5=-1$. \ $f(x)=a(x+1)(x-5)$, $f(2)=-9a=-4{,}5$ $\Rightarrow$ $a=\frac12$:
\[f(x)=\frac12x^2-2x-\frac52\quad(b=-2,\ c=-2{,}5),\qquad n=f(0)=-2{,}5.\]
$g(-1)=l=0$, \ $g(0)=k+l=-2{,}5$ $\Rightarrow$ $k=-2{,}5$.

a) $f(-2)=2+4-2{,}5=3{,}5$.

b) $\dfrac{a-b-l-c}{k-m-n}=\dfrac{0{,}5+2-0+2{,}5}{-2{,}5+1+2{,}5}=\dfrac51=5$.
\Javob{a) $3{,}5$; \quad b) 5}

\newpage
\Savol{39}
a) $f''(0)=4\cdot2-3\cdot1+1=6$.

b) Differensiallaymiz: $f'''=4f''-3f'$ $\Rightarrow$ $f'''(0)=24-6=18$;\\
\phantom{b) }$f^{(4)}=4f'''-3f''$ $\Rightarrow$ $f^{(4)}(0)=72-18=54$.
\Javob{a) 6; \quad b) 54}

\newpage
\Savol{40}
$y=k_2x+l_2$: $(0;6)$, $(12;0)$ $\Rightarrow$ $y=-\frac x2+6$. Perpendikulyar: $k_1=2$, $y=2x+6$.\\
Parabola: $c=6$; $(0;6)$ da urinadi: $f'(0)=b=2$; \ $f(6)=36a+12+6=0$ $\Rightarrow$ $a=-\frac12$:
\[f(x)=-\frac{x^2}2+2x+6.\]
$B$: \ $-\frac{x^2}2+2x+6=-\frac x2+6$ $\Rightarrow$ $x\left(\frac52-\frac x2\right)=0$ $\Rightarrow$ $x=5$, $y=3{,}5$.

a) $5+3{,}5=8{,}5$.

b) Uçburçak $O$, $B(5;3{,}5)$, $(6;0)$: asosi $6$, balandligi $3{,}5$: \ $S=\frac12\cdot6\cdot3{,}5=10{,}5$.
\Javob{a) $8{,}5$; \quad b) $10{,}5$}

\newpage
\Savol{41}
\fig{\begin{tikzpicture}[scale=0.32,font=\small]
\coordinate (A) at (0,0); \coordinate (C) at (14,0); \coordinate (B) at (7,12.124); \coordinate (P) at (5.571,2.227); \coordinate (Q) at (4.714,-3.711);
\fill[okgreen!12] (A)--(B)--(P)--cycle;
\draw (A)--(B)--(C)--cycle; \draw (A)--(P)--(B); \draw (P)--(C);
\draw[axs] (A)--(Q)--(C); \draw[ax] (P)--(Q);
\node[below left] at (A) {$A$}; \node[above] at (B) {$B$}; \node[below right] at (C) {$C$}; \node[above] at (P) {$P$}; \node[aux,below] at (Q) {$P'$};
\node[left] at (2.8,1.1) {$6$}; \node[left] at (6.2,7.2) {$10$};
\end{tikzpicture}}
$OC=R=\dfrac{a}{\sqrt3}=\dfrac{14\sqrt3}3$ $\Rightarrow$ tomon $a=14$.

a) $\triangle ABP$: $6$, $10$, $14$: \ $\cos\angle APB=\dfrac{36+100-196}{120}=-\dfrac12$ $\Rightarrow$ $\angle APB=120^\circ$:
\[S_{ABP}=\frac12\cdot6\cdot10\cdot\sin120^\circ=15\sqrt3.\]
b) $A$ atrofida $60^\circ$ ga buriş ($B\to C$): $P\to P'$. $AP'=AP=6$, $\angle PAP'=60^\circ$ $\Rightarrow$ $APP'$ teng tomonli, $PP'=6$.\\
$P'C=PB=10$, \ $\angle AP'C=\angle APB=120^\circ$ $\Rightarrow$ $\angle PP'C=120^\circ-60^\circ=60^\circ$:
\[PC^2=36+100-2\cdot6\cdot10\cdot\frac12=76\ \Rightarrow\ PC=2\sqrt{19}.\]
\Javob{a) $15\sqrt3$; \quad b) $2\sqrt{19}$}

\newpage
\Savol{42}
\fig{\begin{tikzpicture}[scale=0.22,font=\small]
\coordinate (A) at (0,0); \coordinate (C) at (14,0); \coordinate (B) at (5,12); \coordinate (D) at (5,-12); \coordinate (H) at (5,0); \coordinate (I) at (6.5,0);
\draw (A)--(B)--(C)--(D)--cycle; \draw (A)--(C); \draw (B)--(D); \draw[okgreen!60!black] (I) circle (6);
\fill (I) circle (8pt);
\node[left] at (A) {$A$}; \node[above] at (B) {$B$}; \node[right] at (C) {$C$}; \node[below] at (D) {$D$}; \node[below left] at (H) {$H$}; \node[aux,below right] at (I) {$I$};
\node[above left] at (2.5,6) {$13$}; \node[above right] at (9.5,6) {$15$};
\end{tikzpicture}}
Içki çizilgan törtburçakda $AB+CD=BC+AD$, şuning uçun teng tomonlar qöşni: $AB=AD=13$, $CB=CD=15$ — deltoid. $AC$ — simmetriya öqi, $AC\perp BD$, $BH=HD=h$.\\
$BD$ diagonali $AC$ ni $AH:HC=5:9$ nisbatda böladi ($AH=5k$, $HC=9k$):
\[169-25k^2=225-81k^2\ \Rightarrow\ k=1,\quad AH=5,\ HC=9,\ h=12.\]
a) $S=\frac12\cdot AC\cdot BD=\frac12\cdot14\cdot24=168$.

b) $r=\dfrac Sp=\dfrac{168}{28}=6$. Markaz $I$ — $AC$ da ($A$ burçak bissektrisasi): \ $\sin\angle BAC=\frac{12}{13}$ $\Rightarrow$ $AI=\dfrac r{\sin\angle BAC}=\dfrac{6\cdot13}{12}=6{,}5$.
\[HI=AI-AH=6{,}5-5=1{,}5.\]
\Javob{a) 168; \quad b) $1{,}5$}

\newpage
\Savol{43}
Tomon $a=4$. $AE\perp FC$ ($AE$ — qisqa diagonal, $FC$ — katta diagonal), $FP=\frac a2=2$, $PC=6$, $EP=PA=a\frac{\sqrt3}2=2\sqrt3$.

a) $D$ dan $FC$ gaça masofa $2\sqrt3$: \ $S_{PDC}=\frac12\cdot6\cdot2\sqrt3=6\sqrt3$.

b) $S_{EPF}=\frac12\cdot2\cdot2\sqrt3=2\sqrt3$. \ $ABCP$ — trapetsiya ($AB\parallel PC$), asoslari $4$ va $6$, balandligi $2\sqrt3$: \ $S=\frac{4+6}2\cdot2\sqrt3=10\sqrt3$.
\[\frac{2\sqrt3}{10\sqrt3}=\frac15.\]
\Javob{a) $6\sqrt3$; \quad b) $\frac15$}

\newpage
\Savol{44}
\fig{\begin{tikzpicture}[scale=1.1,font=\small]
\coordinate (P1) at (1,0); \coordinate (P2) at (0.5,0.866); \coordinate (P3) at (-0.5,0.866); \coordinate (P4) at (-1,0); \coordinate (P5) at (-0.5,-0.866); \coordinate (P6) at (0.5,-0.866);
\fill[okgreen!12] (0.634,0.634)--(-0.634,0.634)--(-0.634,-0.634)--(0.634,-0.634)--cycle;
\draw (P1)--(P2)--(P3)--(P4)--(P5)--(P6)--cycle;
\draw[okgreen!60!black,thick] (0.634,0.634)--(-0.634,0.634)--(-0.634,-0.634)--(0.634,-0.634)--cycle;
\draw[ax] (0,0)--(0.634,0.634); \draw[ax] (0,0)--(1,0);
\draw[axs] (0.634,0.634)--(0.634,0); \node[aux,above right] at (0.634,0.634) {$K$}; \node[below] at (0.634,0) {\scriptsize$K_0$}; \node[right] at (P1) {$P$}; \node[below left] at (0,0) {\scriptsize$O$};
\end{tikzpicture}}
Kiçik asos tomoni $s$, katta asos tomoni $2s$ (yuzlar nisbati $4$). Kub yuqori yoği — kiçik asosga içki çizilgan kvadrat (tomon $x$); kvadrat markazi oltiburçak markazida, $K$ uçi oltiburçak tomonida.\\
$K$ dan $OP$ ga perpendikulyar $KK_0=\frac x2$, $OK_0=\frac x2$, $OP=s$; \ $\triangle KK_0P$ da $\angle P=60^\circ$:
\[\tg60^\circ=\frac{x/2}{s-x/2}\ \Rightarrow\ \sqrt3\left(s-\frac x2\right)=\frac x2\ \Rightarrow\ x=\frac{2\sqrt3\,s}{\sqrt3+1}=s\left(3-\sqrt3\right).\]
a) $s=1+\sqrt3$: \ $x=\left(1+\sqrt3\right)\left(3-\sqrt3\right)=2\sqrt3$, \ $V=\left(2\sqrt3\right)^3=24\sqrt3$.

b) $6x^2=24-12\sqrt3$ $\Rightarrow$ $x^2=4-2\sqrt3$, $x=\sqrt3-1$. \ $s=\dfrac{\sqrt3-1}{\sqrt3\left(\sqrt3-1\right)}=\dfrac1{\sqrt3}$, katta asos tomoni $2s=\dfrac{2\sqrt3}3$.
\Javob{a) $24\sqrt3$; \quad b) $\frac{2\sqrt3}3$}

\newpage
\Savol{45}
$A$ va $C$ orasidagi masofaning qirğoqlarga perpendikulyar qismi $30+20+50=100$, parallel qismi $BF$ (ya'ni $BD+EF$):
\[AC^2=100^2+BF^2\ \Rightarrow\ 13\,600=10\,000+BF^2\ \Rightarrow\ BF=60.\]
$DE$ uzunligi özgarmas, $AD+EC$ eng kiçik bölişi kerak. $C$ ni $\overrightarrow{ED}$ ga suramiz: $C\to C'$, $EC=DC'$. $AD+DC'$ eng kiçik, agar $D$ — $AC'$ kesmada, ya'ni $AD\parallel EC$:
\[\triangle ABD\sim\triangle CFE:\ \frac{BD}{EF}=\frac{30}{50}\ \Rightarrow\ BD=\frac{3}{8}\cdot60=22{,}5.\]
b) $AD+EC=AC'=\sqrt{80^2+60^2}=100$, \ $AD+DE+EC=100+20=120$.
\Javob{a) $22{,}5$; \quad b) 120}

\ShowJavoblar
\end{document}
